\newpage
\section{Заключение}

В работе разработана и реализована система AI Coauthor — комплекс для семантического анализа
научного пространства, тематической кластеризации публикаций и рекомендации потенциальных
соавторов с объяснениями на основе LLM. Система объединяет офлайн-предподсчёт эмбеддингов,
кластеризацию, построение метакластеров, обучение retrieval-моделей и онлайн-этап
LLM-переранжирования.

Эмпирическая часть построена на корпусе из 509\,447 публикаций OpenAlex по подполю Artificial
Intelligence. Статьи представлены 768-мерными эмбеддингами multilingual-e5-base. Для карты
научного пространства выбран BERTopic-style pipeline: UMAP, HDBSCAN, c-TF-IDF, MMR и
LLM-разметка тем. Итоговая структура содержит 979 fine-тем и 76 метакластеров. HDBSCAN выбран по
совокупности метрик качества тем, а не только по внутренним геометрическим критериям.

В рекомендательной части сравнивались mean-бейзлайн, Transformer-представление автора, GraphSAGE
по графу соавторства и гибрид Transformer+GraphSAGE. Финальный LLM-этап работает с top-10
кандидатами retrieval и выбирает 3 рекомендации. Качество финальной выдачи подтверждено
LLM-оценкой: LLM Selected3Mean@10 достигает 2,492 из 3, а LLM Has3@10 при score $\geq 2$ равен
0,875. Observed-метрики по будущим соавторам используются как строгий proxy retrieval, но не как
основное подтверждение релевантности рекомендаций.

Полученные результаты показывают, что граф соавторства даёт качество, сопоставимое с
Transformer-представлением автора, но явного значимого выигрыша от метакластерных узлов в retrieval
не получено. Поэтому кластеризация в текущей версии системы прежде всего полезна как слой
интерпретации: она делает карту научного пространства читаемой, даёт названия тем и связывает
рекомендации с тематической структурой корпуса. Её вклад в качество извлечения кандидатов требует
дополнительной проверки на более широкой выборке и более строгой LLM-evaluation.

Основные ограничения связаны с жёстким назначением статей в темы, искажениями 2D-проекции,
разреженностью будущих соавторств как target-разметки и тем, что в работе используются только
заголовки и аннотации, без полного текста публикаций. Направления развития — мягкое overlapping
назначение тем, использование полного текста, расширение корпуса за пределы Artificial
Intelligence и автоматизация полного воспроизводимого пайплайна.
